% Part: computability
% Chapter: computability-theory
% Section: complete-ce-sets

\documentclass[../../../include/open-logic-section]{subfiles}

\begin{document}

\olfileid{cmp}{thy}{cce}
\olsection{完全可计算枚举集}

\begin{defn}
一个集合~$A$ 是\emph{完全可计算枚举集}（在多一可归约意义下），如果
\begin{enumerate}
\item $A$ 是可计算枚举的，并且
\item 对于任何其他可计算枚举集~$B$，都有 $B \leq_m A$。
\end{enumerate}
\end{defn}

换言之，完全可计算枚举集是所有可能的可计算枚举集中“最难”的。借助它们，能够回答关于\emph{任何}可计算枚举集的问题。

\begin{thm}
$K$、$K_0$ 和 $K_1$ 都是完全可计算枚举集。
\end{thm}

\begin{proof}
为证 $K_0$ 是完全的，设 $B$ 是任意可计算枚举集。则对某个指标 $e$，\[
B = W_e = \Setabs{x}{\cfind{e}(x) \fdefined}.
\]
令 $f$ 为函数 $f(x) = \tuple{e, x}$。那么对每个自然数 $x$，$x \in B$ 当且仅当 $f(x) \in K_0$。换言之，$f$ 将 $B$ 归约到~$K_0$。

为证 $K_1$ 是完全的，注意到在 \olref[k1]{prop:k1} 的证明中我们已将 $K_0$ 归约到了 $K_1$。因此，根据 \olref[ppr]{prop:trans-red}，任何可计算枚举集也都能归约到~$K_1$。

$K$ 可以用几乎同样的方式归约到 $K_0$。
\end{proof}

\begin{prob}
给出 $K$ 到 $K_0$ 的一个归约。
\end{prob}

\begin{digress}
这样一来，我们迄今考虑过的所有可计算枚举集的例子，要么是可计算的，要么是完全的。这似乎有些奇怪！是否存在既不可计算也不完全的可计算枚举集的例子呢？答案是肯定的，但直到20世纪50年代中期，这一点才由 Friedberg 和 Muchnik 独立证明。
\end{digress}

\end{document}